\section{Auswertung}


\subsection{2. Aufgabe}
In dieser Aufgabe werden nun die Ergebnisse aus Teil 2 und 3 des Experiments ausgewertet. In dieser Aufgabe sind jeweils die Frequenzen $\omega_1$ und $\omega_2$ der Normalschwingungen sowie $\omega_3$ und $\omega_s$ der Schwebung für beide Kopplungsgrade zu bestimmen. Das Vorgehen ist dabei analog zu Aufgabe 1. Die gemessenen Zeiten entsprechen der Zeit von 8 Schwingungen. 

\subsection*{Leichtes Gewicht 1. und 2. Normalschwingung}
In diesem Abschnitt werden die Frequenzen $\omega_1$ und $\omega_2$ der 1. und 2. Normalschwingung bestimmt. Mit der obigen Formel ergibt sich folgende Tabelle für die einzelnen Werte. 
\begin{table}[H]
    \centering
    \begin{tabular}{c|c}
        1. Normalschwingung & 2. Normalschwingung \\
        \hline
        $\omega_1 \left(\text{Hz}\right)$ & $\omega_2 \left(\text{Hz}\right)$ \\
        \hline
        $3{,}61 \pm 0{,}026$ & $3{,}75 \pm 0{,}028$ \\
        $3{,}61 \pm 0{,}026$ & $3{,}69 \pm 0{,}027$ \\
        $3{,}72 \pm 0{,}027$ & $3{,}72 \pm 0{,}027$ \\
        $3{,}61 \pm 0{,}026$ & $3{,}72 \pm 0{,}027$ \\
        $3{,}64 \pm 0{,}026$ & $3{,}75 \pm 0{,}028$ \\
    \end{tabular}
    \caption{Die ermittelten Frequenzwerte mit den dazugehörigen Fehlern von $\omega_1$ und $\omega_2$}
\end{table}
Die Mittelwerte der Tabelle ergeben
\begin{equation*}
    \centering
    \bar{\omega}_1 \approx \left(3{,}63 \pm 0{,}0262 \right)\,\text{Hz}
\end{equation*}
\begin{equation*}
    \centering
    \bar{\omega}_2 \approx \left(3{,}73 \pm 0{,}0274 \right)\,\text{Hz}
\end{equation*}
Die Ergebnisse mit einem relativen Fehler von circa 1\,\% wirken zuverlässig. Der Unterschied zwischen den beiden Mittelwerten kommt daher, dass die Schwingungsperiode der 2. Normalschwingung stärker vom Kopplungsgrad abhängt.

\subsection*{Leichtes Gewicht Schwebung}
In diesem Teil werden $\omega_3$ und $\omega_s$ für das leichte Kopplungsgewicht ermittelt. Bei $\omega_s$ muss eine andere Anzahl an Schwingungen beachtet werden, da die einhüllende Schwingung um einiges langsamer ist als die normalen Pendelschwingungen. 
\begin{table}[H]
    \centering
    \begin{tabular}{c|c}
        $\omega_3 \left(\text{Hz}\right)$ & $\omega_s \left(\text{Hz}\right)$ \\
        \hline
        $3{,}81 \pm 0{,}029$ & $0{,}098 \pm 0{,}001$ \\
        $3{,}72 \pm 0{,}027$ & $0{,}099 \pm 0{,}001$ \\
        $3{,}81 \pm 0{,}029$ & $0{,}098 \pm 0{,}001$ \\
        $3{,}81 \pm 0{,}029$ & $0{,}098 \pm 0{,}001$ \\
        $3{,}84 \pm 0{,}029$ & $0{,}099 \pm 0{,}001$ \\
    \end{tabular}
    \caption{Die ermittelten Frequenzwerte mit den dazugehörigen Fehlern von $\omega_3$ und $\omega_s$}
\end{table}
Die Mittelwerte betragen somit
\begin{equation*}
    \centering
    \bar{\omega}_3 \approx \left(3{,}798 \pm 0{,}0286 \right)\,\text{Hz}
\end{equation*}
\begin{equation*}
    \centering
    \bar{\omega}_s \approx \left(0{,}0984 \pm 0{,}001 \right)\,\text{Hz}
\end{equation*}
Wie bereits beim Versuch festgestellt wurde, ist die einhüllende Frequenz der Schwebung deutlich kleiner als die der einzelnen Schwingungen. Mit einem relativen Fehler von circa 1\,\% scheinen die Ergebnisse zuverlässig.

\subsection*{Schweres Gewicht 1. und 2. Normalschwingung}
Nun wird der Kopplungsgrad des schweren Gewichtes betrachtet, auch hier werden wieder die Ergebnisse der 1. und 2. Normalschwingung ausgewertet. Aus den Messdaten ergibt sich folgende Tabelle.
\begin{table}[H]
    \centering
    \begin{tabular}{c|c}
        $\omega_1 \left(\text{Hz}\right)$ & $\omega_2 \left(\text{Hz}\right)$ \\
        \hline
        $3{,}61 \pm 0{,}026$ & $4{,}33 \pm 0{,}037$ \\
        $3{,}59 \pm 0{,}026$ & $4{,}15 \pm 0{,}034$ \\
        $3{,}72 \pm 0{,}027$ & $4{,}22 \pm 0{,}035$ \\
        $3{,}72 \pm 0{,}027$ & $4{,}09 \pm 0{,}039$ \\
        $3{,}69 \pm 0{,}027$ & $4{,}09 \pm 0{,}039$ \\
    \end{tabular}
    \caption{Die ermittelten Frequenzwerte mit den dazugehörigen Fehlern von $\omega_1$ und $\omega_2$}
\end{table}
Die Mittelwerte ergeben
\begin{equation*}
    \centering
    \bar{\omega}_1 \approx \left(3{,}666 \pm 0{,}0266 \right)\,\text{Hz},
\end{equation*}
\begin{equation*}
    \centering
    \bar{\omega}_2 \approx \left(4{,}176 \pm 0{,}184 \right)\,\text{Hz}.
\end{equation*}
Ähnlich wie beim ersten Teil unterscheiden sich die Frequenzen der 1. und 2. Normalschwingung, diesmal ist der Unterschied jedoch deutlich ausgeprägter. Damit stimmen Theorie und Praxis in diesem Punkt überein.

\subsection*{Schweres Gewicht Schwebung}
In diesem Teil werden $\omega_3$ und $\omega_s$ für das schwere Kopplungsgewicht ermittelt. Auch hier muss bei $\omega_s$ die Anzahl der Schwingungen beachtet werden, da die einhüllende Schwingung deutlich langsamer ist als die normalen Pendelschwingungen.
\begin{table}[H]
    \centering
    \begin{tabular}{c|c}
        $\omega_3 \left(\text{Hz}\right)$ & $\omega_s \left(\text{Hz}\right)$ \\
        \hline
        $3{,}69 \pm 0{,}027$ & $0{,}490 \pm 0{,}0002$ \\
        $3{,}69 \pm 0{,}027$ & $0{,}499 \pm 0{,}0002$ \\
        $3{,}69 \pm 0{,}027$ & $0{,}491 \pm 0{,}0002$ \\
        $3{,}69 \pm 0{,}027$ & $0{,}504 \pm 0{,}0002$ \\
        $3{,}75 \pm 0{,}028$ & $0{,}504 \pm 0{,}0002$ \\
    \end{tabular}
    \caption{Die ermittelten Frequenzwerte mit den dazugehörigen Fehlern von $\omega_3$ und $\omega_s$}
\end{table}
Die Mittelwerte betragen
\begin{equation*}
    \centering
    \bar{\omega}_3 \approx \left(3{,}702 \pm 0{,}0272 \right)\,\text{Hz}
\end{equation*}
\begin{equation*}
    \centering
    \bar{\omega}_s \approx \left(0{,}497 \pm 0{,}0002 \right)\,\text{Hz}
\end{equation*}
Die Werte für $\omega_3$ und $\omega_s$ unterscheiden sich deutlich, was ebenfalls mit der Theorie vereinbar ist. Beim Vergleich der Werte des leichten und des schweren Kopplungsgewichtes ist eine klare Abhängigkeit zwischen den Ergebnissen erkennbar, auch dies stimmt mit der Theorie überein.

\subsection*{Fehlerbetrachtung}
Wie in Aufgabe 1 sind hier die gleichen Fehlerquellen zu beachten. Hauptsächliche Fehler sind die tatsächlichen Zeitpunkte im Gegensatz zu den gemessenen Zeiten, da beim Bedienen der Stoppuhr leicht Zeitabweichungen entstehen können. Ein weiterer relevanter Fehler sind die Reibungskräfte, die in der Realität auftreten, in der Theorie jedoch vernachlässigt wurden. Des Weiteren wurde in der Theorie stets mit der Kleinwinkelnäherung gerechnet, obwohl es leicht passieren kann, dass das Pendel über $5^\circ$ ausgelenkt wird. Als letzter Punkt ist anzumerken, dass der Zeitpunkt des Loslassens nicht perfekt gleichzeitig war, bedingt durch menschliche Fehler.

